\section{Security Analysis}
\label{sec:security}

We argue each goal of \S\ref{sec:threat} from the rules. Signatures are
EUF-CMA Ed25519, and SHA-256 is collision-resistant.

\paragraph{G1 Attributability.}
By R-1, a counted signature is a valid signature by the key named in the
core. By R-4, positive evidence requires the requester's signature.
Forging positive evidence against an honest requester therefore requires
forging that requester's signature. A requester-only receipt is negative
evidence, and it is attributable to its signer, whose weight is the
observer's decision (T-3). Role tags in the signing input stop a receipt
signature from being replayed as a checkpoint signature, or the reverse.
Verifying the received bytes rules out re-encoding tricks.

\paragraph{G2 Local soundness.}
Under T-3 an issuer's weight is 1 for the observer itself, and otherwise
a function of evidence the observer verified. Under W1 it is agreement
between the issuer's ratings and the observer's own. An issuer that
shares fewer than three rated providers with the observer gets weight 0.
A Sybil ring that only rates itself therefore contributes nothing. The
rebuild rule (T-1/T-2) makes the projection a pure function of the
receipt store, so nothing, including an implementation bug on one
ingestion path, can raise a score without evidence.

What G2 does \emph{not} give is Sybil-proofness against an adversary who
buys agreement with real work. Camouflaging Sybils earn W1 weight by
rating honest providers truthfully. W1 bounds what \emph{one} Sybil buys
to its direct agreement with the observer, but not what many buy
together. W1 falls behind self-only weighting once camouflaging Sybils
outnumber honest requesters about six to one (A10). W2's propagation
fails sooner (\S\ref{sec:eval-trust}). The weight that camouflage buys can
also be aimed. Camouflaged Sybils that slander five chosen honest
providers cut their traffic by 68\% under W1 and 88\% under W2 (A12). Cheng and Friedman~\cite{cheng2005sybilproof}
showed that no
\emph{symmetric} reputation function is Sybil-proof. \sys{}'s weighting
is asymmetric, because each observer computes its own. That does not
help against camouflage: those Sybils earn their agreement with real,
truthful tasks, so evidence alone cannot tell them from honest raters. It
can still bound the damage. DSybil~\cite{yu2009dsybil} bounds each user's
loss against unlimited Sybils, under an assumption about how honest users
vote. W1 carries no such bound, and adapting a loss-bounded weighting to
receipts is future work. \sys{}'s answer for now is to make evidence cost
real tasks, to limit what agreement buys to direct agreement (W1), and to
recommend a canary in front.

\paragraph{G3 Safe resumption.}
A valid resumption point carries both signatures, and the requester
signs only after recomputing the root (C-2). A successor recomputes the
root over the replayed prefix and refuses to continue on a mismatch. So a
successor never continues from content that the requester did not verify,
unless it finds a second preimage for the root. By C-5, a failed
provider's claim is bounded by the usage in the last checkpoint that
both parties signed. C-3/C-4 stop a provider from splicing an older,
higher-usage checkpoint into the chain. The zone label is inside the
signed core, so a successor cannot be handed the task under a wider zone
without the requester signing that wider label (Z-3).

\paragraph{G4 Bounded exposure.}
Suppose vouchers are issued per checkpoint and the requester follows M-6.
Then the requester never authorizes more than the price of acknowledged
usage plus the credit it chooses to extend. That is not one interval: its
exposure is whatever credit it decides to extend. The provider never
delivers more than one interval ahead of a covering voucher. Its unpaid
work per task is therefore at most $\min(L,c)$ Parts on the fallback
path, plus the 2.5--5 Parts in flight on the primary path
(\S\ref{sec:eval-meter}). The simulation checks both bounds in every run
(H5), but they hold by construction.

\paragraph{Dishonest requesters.}
\sys{} protects requesters from providers far better than it protects
providers from requesters.
\begin{itemize}
\item \textbf{Free-riding.} A requester that stops countersigning or
  paying gets at most the unpaid work above for free: $\min(L,c)$ Parts
  per task. Requester identities are free, and \sys{} computes no trust
  \emph{in requesters}, so a free-rider can repeat this indefinitely.
  The C-6 trade-off sharpens it: at the synchronous default $c{=}32$, any
  task of 32 Parts or fewer can be taken whole. A provider's defences are
  local policy, outside this layer: a short first interval, a rail that
  takes payment up front, or its own record of requesters.
\item \textbf{Receipt refusal.} A requester can refuse to countersign the
  final receipt. That withholds positive evidence the provider earned,
  and R-4 leaves the provider no way to record it alone.
\item \textbf{Vindictive customers.} A genuine customer's unilateral
  negative receipt is admitted, weighted by the observer's agreement with
  that customer. The provider can dispute it but cannot disprove it,
  which is the same asymmetry \S\ref{sec:intro} criticises in
  provider-signed receipts, with the roles reversed. Nothing binds the
  receipt to a real task either, because the provider chose the A2A task
  identifier and never signed. Issuer weighting bounds the damage one
  customer can do (T-3, the per-issuer cap). It does not bound the damage
  a customer the observer already agrees with can do.
\end{itemize}

\paragraph{Reviewer findings.}
Before this work, a role-based review demonstrated exploits against the
first version of the code. It ran as eight AI-model sessions with no
access to the implementation sessions (Acknowledgments), and was not a
human review. Here is how each finding was resolved:
\begin{itemize}
\item \textbf{Receipt farming.} A Sybil ring reached a score of 0.998 in
  about 150\,ms of receipt generation. Fixed by R-2 to R-4, T-2 and T-3.
  E1 measures T-3's part. R-4's part is the G1 argument above, because
  no E1 agent emits a receipt that R-4 rejects.
\item \textbf{Shard flooding.} Fixed by D-6; measured in E3a.
\item \textbf{Negative amounts.} Made unencodable.
\item \textbf{Unsigned checkpoint usage.} Fixed by the new checkpoint
  artifact.
\item \textbf{Replay-cache cost and TTL.} Fixed.
\item \textbf{Identifier grinding.} Not fixed. The proposed mitigation
  was measured and found ineffective (E3b), and the specification now
  lists grinding as a limitation. Captured shards contain identities with
  no evidence, but the trust layer can rank only what discovery returns.
  Once discovery is flooded, it cannot help unless requesters keep
  known-good providers on their shortlists, and new requesters have none
  (A11, \S\ref{sec:eval-discovery}).
\end{itemize}
